\section{Discusión y retos actuales}

La incorporación de emociones dentro de agentes inteligentes ha permitido desarrollar sistemas más adaptativos, creíbles y socialmente coherentes. A lo largo de los últimos años, la computación afectiva ha evolucionado desde modelos emocionales simples hasta arquitecturas complejas capaces de integrar razonamiento cognitivo, emociones y comportamiento social.

Sin embargo, su representación computacional continúa presentando importantes limitaciones. La fuerte dependencia contextual, cultural y subjetiva del afecto humano resulta difícil de modelar mediante representaciones simplificadas. De hecho, muchos sistemas actuales utilizan conjuntos limitados de emociones o mecanismos estáticos, lo que reduce la naturalidad del agente. A esto se suma la complejidad de la validación experimental, ya que evaluar si un agente expresa emociones de manera realista y coherente sigue siendo un problema abierto en la Inteligencia Artificial social.

En general, las líneas de investigación actuales buscan desarrollar agentes capaces de mantener estados emocionales más coherentes a largo plazo, integrar aprendizaje y memoria afectiva, y adaptarse de manera flexible a distintos contextos sociales y culturales.

\subsection{Relación entre los trabajos analizados}

Los trabajos estudiados representan diferentes etapas dentro de la evolución de los agentes afectivos. El modelo OCC proporciona una estructura cognitiva para generar emociones mediante procesos de appraisal, mientras que arquitecturas como EBDI integran dichas emociones dentro del razonamiento deliberativo de agentes BDI.

Por otro lado, WASABI extiende estos enfoques hacia humanos virtuales capaces de representar dinámicas emocionales continuas mediante el modelo PAD. Asimismo, los trabajos de Taverner et al. incorporan representación difusa y adaptación cultural, resolviendo algunas limitaciones presentes en modelos emocionales puramente discretos.

Desde nuestro punto de vista, cada enfoque presenta ventajas y limitaciones diferentes. OCC resulta especialmente útil por su interpretabilidad semántica y su facilidad para relacionar emociones con procesos de appraisal, aunque los modelos discretos presentan dificultades para representar intensidad y evolución temporal de manera continua. Por otro lado, modelos dimensionales como PAD permiten describir transiciones emocionales más suaves, pero sacrifican parte de la interpretación semántica directa de cada emoción. Arquitecturas como EBDI integran emociones dentro de la toma de decisiones, mientras que WASABI destaca especialmente en contextos orientados a interacción y expresividad emocional.

Finalmente, los agentes normativos emocionales amplían estas arquitecturas hacia entornos sociales y multiagente, donde emociones como culpa, empatía o vergüenza influyen directamente sobre el comportamiento cooperativo y el cumplimiento de normas.

En definitiva, la computación afectiva está adquiriendo un papel cada vez más relevante dentro de la Inteligencia Artificial social.